\\newpage
\\section{误差分析}

\\subsection{模型误差来源}
\\begin{enumerate}
    \\item[(1)] \\textbf{参数估计误差}：MNL模型中的价格敏感系数$\\beta，存在统计误差。
    \\item[(2)] \\textbf{假设简化误差}：假设用户仅考虑价格和距离，忽略了其他因素（如停车场设施、安全性等）。
    \\item[(3)] \\textbf{数据噪声}：原始数据存在测量误差和缺失值。
\\end{enumerate}

\\subsection{灵敏度分析}
分析价格弹性系数$\\beta：
\\begin{table}[H]
\\centering
\\caption{价格弹性灵敏度分析}
\\begin{tabular}{c c c}
\\toprule
$\\beta & 最优价格变化 & 收益变化 \\\\
\\midrule
-20\\% & -3.2\\% & +1.5\\% \\\\
-10\\% & -1.8\\% & +0.8\\% \\\\
基准 & 0 & 0 \\\\
+10\\% & +2.1\\% & -0.5\\% \\\\
+20\\% & +4.5\\% & -2.1\\% \\\\
\\bottomrule
\\end{tabular}
\\end{table}

结果表明，当$\\beta\\%时，最优价格上升约2\\%，总收益下降约0.5\\%，模型具有较好的鲁棒性。
